\clearpage
\VerdierPhysicalPageAnchor{0185}%
\begingroup\small\setlength{\parskip}{0pt}\setlength{\emergencystretch}{1.5em}
\setlength{\abovedisplayskip}{4pt plus 2pt minus 2pt}%
\setlength{\belowdisplayskip}{4pt plus 2pt minus 2pt}%

\VerdierSemanticAnchor{ch3d-1-3-6-conditions-cont}
\noindent\textit{i) $F$ 通过 $K(\mathcal A)$ 分解。}

\noindent\textit{ii) $F$ 把同伦等价变为 $C$ 中的同构。}

\VerdierSemanticAnchor{ch3d-1-3-6-proof}
i) $\Rightarrow$ ii) 是显然的。为证明 ii) $\Rightarrow$ i)，只须证明
$F$ 把每对相互同伦的态射变为相等的态射。设 $X$ 是
$\operatorname{comp}(\mathcal A)$ 的一个对象。$\operatorname{comp}(\mathcal A)$
上的函子
\[
Y\longmapsto
\{\text{从 $X$ 到 $Y$ 的态射对，配有一个将二者相连的同伦}\}
\]
由复形 $Z$ 表示：

\VerdierSemanticAnchor{ch3d-1-3-6-representing-complex}
\begin{equation}
\begin{aligned}
Z^i&=X^i\oplus X^i\oplus X^{i+1},\\[2pt]
d_Z^i&=
\begin{pmatrix}
d_X^i&0&\operatorname{id}_{X^{i+1}}\\
0&d_X^i&-\operatorname{id}_{X^{i+1}}\\
0&0&-d_X^{i+1}
\end{pmatrix}.
\end{aligned}
\tag{1.3.6.1}
\end{equation}
因而有两个相互同伦的态射 $p_1,p_2\colon X\longrightarrow Z$：

\VerdierSemanticAnchor{ch3d-1-3-6-morphisms}
\begin{equation}
p_1^i=
\begin{pmatrix}
\operatorname{id}_{X^i}\\0\\0
\end{pmatrix},
\qquad
p_2^i=
\begin{pmatrix}
0\\\operatorname{id}_{X^i}\\0
\end{pmatrix},
\tag{1.3.6.2}
\end{equation}
只须证明 $F(p_1)=F(p_2)$。而态射对
$(\operatorname{id}_X,\operatorname{id}_X)$ 连同平凡同伦定义一个态射
$\pi\colon Z\longrightarrow X$：

\VerdierSemanticAnchor{ch3d-1-3-6-retraction}
\begin{equation}
\pi^i=(\operatorname{id}_{X^i},\operatorname{id}_{X^i},0)
\tag{1.3.6.3}
\end{equation}
并且有 $\pi p_1=\pi p_2=\operatorname{id}_X$。此外，立即验证态射
$\pi$、$p_1$ 和 $p_2$ 都是同伦等价。因此，$F(p_1)$、$F(p_2)$ 和
$F(\pi)$ 都是同构，并且
$F(\pi)F(p_1)=F(\pi)F(p_2)=\operatorname{id}_{F(X)}$。于是
$F(p_1)=F(p_2)$，引理证毕。

\VerdierSemanticAnchor{ch3d-1-3-7}
\medskip\noindent\textbf{注 1.3.7.}\quad
由命题 (1.3.5)，$\mathcal A$ 的导出范畴是如下普遍问题的解：在
$\operatorname{comp}(\mathcal A)$ 中逆转拟同构（第二章，2.2.3，b)）。
\endgroup
